\chapter{Thanh ghi đếm thích nghi cho
  \texorpdfstring{ràng buộc đúng \(K\)}{ràng buộc đúng K}}
\chapterlead{Bộ đếm tuần tự thông thường dành cùng một độ rộng cho mọi
hàng. Thanh ghi thích nghi chỉ tạo những trạng thái đếm có thể đạt tới, nhờ
đó giữ được cấu trúc tuần tự nhưng loại bỏ phần tam giác trạng thái bất khả.}

\index{exactly-k}\index{New Sequential Counter}
\index{adaptive counter}\index{Social Golfer}

\flowdiagram{Dãy Boolean}{Miền trạng thái khả đạt}{Truy hồi ngưỡng}
  {Điều kiện biên}{Đúng \(K\)}

\section{Ràng buộc đúng \texorpdfstring{\(K\)}{K} trong không gian thiết kế}

\emph{Exactly-\(K\)} (ràng buộc đúng \(K\)) có nhiều cách phân rã đúng:
\[
  \sum X=K
  \iff
  \left(\sum X\le K\right)\land
  \left(\sum X\ge K\right),
\]
nhưng hai vế không nhất thiết phải dùng hai bộ biểu diễn độc lập. Một vectơ đơn
phân đầy đủ cho phép đọc cùng lúc \(o_K\) và \(\neg o_{K+1}\); BDD giữ các
trạng thái tổng; mạng sắp xếp các đầu vào; bộ đếm tuần tự theo dõi tiền tố. Câu
hỏi của chương là cụ thể: nếu đã chọn các trạng thái ngưỡng tuần tự thì những
ô nào thật sự có thể đạt và cần được hiện thực hóa (reification)?

\begin{center}
\captionof{table}{Các hướng biểu diễn ràng buộc đúng \(K\).}
\label{tab:exactly-k-designs}
\begin{tabularx}{\textwidth}{@{}>{\bfseries\sffamily}p{.19\textwidth}YYY@{}}
\toprule
Họ & Trạng thái & Giao diện EXK & Đặc trưng\\
\midrule
Hai AMK & Hai bộ đếm & Hai cận tách & Dễ ghép nhưng có thể lặp trạng thái\\
Tuần tự & Tổng tiền tố & Chống tràn + đạt \(K\) & Tốt khi \(K\) nhỏ\\
Cây tổng & Tổng cây con & \(o_K\land\neg o_{K+1}\) & Dùng lại nhiều cận\\
Mạng & Vị trí đã sắp & Hai đầu ra kề nhau & Cấu trúc đều, độ sâu thấp\\
BDD & \((i,s)\) & Chấp nhận trạng thái \(K\) & Hợp nhất trạng thái theo tổng\\
NSC & Ngưỡng tiền tố khả đạt & Biên trên + dưới & Cắt vùng tam giác bất khả\\
\bottomrule
\end{tabularx}
\end{center}

Các cận dưới cho AMO và những khảo cứu về ràng buộc đếm đặt số biến bên cạnh
số mệnh đề, literal và mức lan truyền
\parencite{kucera2019lower,bailleux2010survey}. NSC giữ bậc \(O(nK)\), giảm
một phần xác định của bảng và thay cách đóng biên Exactly-\(K\). Đóng góp cốt
lõi là một \emph{không gian trạng thái gọn hơn trong họ tuần tự}; vị trí trên
biên Pareto (Pareto frontier) được đo bằng số mệnh đề, số literal và sức lan truyền.

\section{Từ bảng chữ nhật đến miền trạng thái khả đạt}

Xét dãy Boolean \(X=\langle x_1,\ldots,x_n\rangle\) và ràng buộc
\[
  \sum_{i=1}^{n}x_i=k.
\]
Một bộ đếm tuần tự có thể dành \(k\) cột cho mọi tiền tố. Tuy nhiên, sau
\(i\) biến đầu tiên, tổng không thể lớn hơn \(i\). Các ô \(j>i\) vì thế không
mang trạng thái hợp lệ.

\emph{New Sequential Counter} (bộ đếm tuần tự mới, NSC) dùng thanh ghi tam giác
\[
  R_{i,j}=1
  \quad\Longleftrightarrow\quad
  \sum_{q=1}^{i}x_q\ge j,
  \qquad
  \begin{matrix}
    1\le i\le n-1,\\[-1mm]
    1\le j\le \min(i,k).
  \end{matrix}
\]
Chỉ \(n-1\) tiền tố được hiện thực hóa (reification); biến \(x_n\) được xử lý trong khi giải (in-processing) khi giải (in-processing) điều kiện cuối.
Với \(n>k\), số biến phụ là
\[
  \sum_{i=1}^{n-1}\min(i,k)
  =
  k(n-1)-\frac{k(k-1)}{2}.
\]
So với bảng chữ nhật \(k(n-1)\), phần tiết kiệm chính xác là
\(k(k-1)/2\) biến. Bậc tiệm cận được giữ nguyên, còn hệ số hữu hạn giảm đúng các
trạng thái bất khả mà bộ giải phải duy trì.

\begin{figure}[tb]
\centering
\begin{tikzpicture}[satfig,x=7mm,y=7mm]
  \node[satlabel] at (2.5,1.25) {bảng chữ nhật};
  \node[satlabel] at (9.5,1.25) {miền khả đạt của NSC};
  \foreach \i in {1,...,5}{
    \foreach \j in {1,...,3}{
      \draw[Rule,fill=SoftGray] (\i,-\j) rectangle ++(.82,.82);
      \node[font=\sffamily\tiny] at (\i+.41,-\j+.41) {\(R_{\i,\j}\)};
      \ifnum\j>\i
        \fill[Gold!28] (\i,-\j) rectangle ++(.82,.82);
        \draw[Gold,thick] (\i,-\j)--++(.82,.82);
        \draw[Gold,thick] (\i+.82,-\j)--++(-.82,.82);
      \fi
    }
  }
  \foreach \i in {1,...,5}{
    \foreach \j in {1,...,3}{
      \ifnum\j>\i
      \else
        \draw[Blue,fill=PaleBlue] ({7+\i},-\j) rectangle ++(.82,.82);
        \node[font=\sffamily\tiny] at ({7+\i+.41},-\j+.41) {\(R_{\i,\j}\)};
      \fi
    }
  }
  \node[satlabel,text=Gold,align=center] at (2.8,-4)
    {các ô \(j>i\)\\không thể đạt};
  \node[satlabel,text=Blue,align=center] at (10.1,-4)
    {\(12\) ô thay vì \(15\)\\khi \(n=6,k=3\)};
\end{tikzpicture}
\caption{Bộ đếm tuần tự hình chữ nhật và miền trạng thái tam giác. NSC không
tạo ba ô bất khả \(R_{1,2},R_{1,3},R_{2,3}\).}
\label{fig:nsc-triangular-domain}
\end{figure}

\begin{keyidea}{``Thích nghi'' ở đây có nghĩa gì?}
Độ rộng hàng \(i\) là \(\min(i,k)\), được quyết định bởi số trạng thái có thể
đạt sau \(i\) đầu vào. ``Thích nghi'' ở đây mô tả hình dạng tĩnh của bộ đếm,
được xác định trước khi bộ giải chạy.
\end{keyidea}

\section{Đặc tả ngữ nghĩa và hệ thức truy hồi}

Dùng hai quy ước biên
\[
  R_{i,0}\equiv\top,
  \qquad
  R_{0,j}\equiv\bot\quad(j\ge1).
\]
Khi đó mọi ô hợp lệ tuân theo
\begin{equation}
  R_{i,j}
  \;\longleftrightarrow\;
  R_{i-1,j}\lor
  \bigl(x_i\land R_{i-1,j-1}\bigr).
  \label{eq:nsc-recurrence}
\end{equation}
Vế phải mô tả đúng hai cách để tiền tố mới đạt ngưỡng \(j\): tiền tố cũ đã đạt
ngưỡng, hoặc \(x_i=1\) và tiền tố cũ đã đạt \(j-1\).

Trong CNF, hệ thức truy hồi được phân rã thành các quan hệ kéo theo theo hai chiều. Phần
tiến gồm
\[
  R_{i-1,j}\rightarrow R_{i,j},
  \qquad
  (x_i\land R_{i-1,j-1})\rightarrow R_{i,j}.
\]
Phần lùi gồm
\[
  (\neg x_i\land\neg R_{i-1,j})\rightarrow\neg R_{i,j},
  \qquad
  \neg R_{i-1,j-1}\rightarrow\neg R_{i,j}.
\]
Sau khi thay các hằng biên và bỏ literal thừa, ta thu được các mệnh đề hai hoặc
ba literal. Cách sinh mệnh đề trực tiếp theo miền tam giác tránh tạo biến rồi
giao cho bộ tiền xử lý xóa đi.

\begin{designrule}{Thanh ghi được đọc như một đại lượng ngữ nghĩa}
Nếu \(R_{i,j}\) xuất hiện trong điều kiện biên, hàm mục tiêu hoặc một ràng buộc
khác, hệ thức truy hồi phải bảo đảm đặc tả hai chiều. Chỉ các quan hệ kéo theo
tiến không đủ để diễn giải \(\neg R_{i,j}\) là ``tiền tố chưa đạt \(j\)''.
\end{designrule}

\section{Đóng ràng buộc \texorpdfstring{đúng \(K\)}{đúng K}}

Phần nhiều nhất \(k\) cấm một đầu vào mới làm tràn bộ đếm:
\begin{equation}
  x_i\rightarrow\neg R_{i-1,k},
  \qquad i=k+1,\ldots,n.
  \label{eq:nsc-overflow}
\end{equation}
Phần ít nhất \(k\) yêu cầu tổng cuối đạt ngưỡng \(k\):
\begin{equation}
  R_{n-1,k}\lor
  \bigl(x_n\land R_{n-1,k-1}\bigr).
  \label{eq:nsc-lower}
\end{equation}
Biểu thức \cref{eq:nsc-lower} được phân phối thành CNF trước khi xuất. Khi
\(k=1\), \(R_{n-1,0}\equiv\top\); khi \(k=n\), có thể đơn giản hóa toàn bộ
ràng buộc thành các mệnh đề đơn vị \(x_i\).

\begin{theorem}[Đúng đắn của NSC]
Hệ thức truy hồi \cref{eq:nsc-recurrence} cùng các điều kiện
\cref{eq:nsc-overflow,eq:nsc-lower} tương đương theo phép chiếu với
\(\sum_{i=1}^{n}x_i=k\).
\end{theorem}

\begin{proof}
Hệ thức truy hồi và các hằng biên cho, bằng quy nạp theo \(i\),
\[
  R_{i,j}=1
  \Longleftrightarrow
  \sum_{q=1}^{i}x_q\ge j.
\]
Nếu có hơn \(k\) đầu vào đúng, xét đầu vào đúng thứ \(k+1\), ký hiệu \(x_i\).
Ngay tại vị trí trước đó, \(R_{i-1,k}=1\), trái với
\cref{eq:nsc-overflow}. Do đó tổng không vượt \(k\).

Ngược lại, \cref{eq:nsc-lower} chỉ đúng khi \(k\) đầu vào đầu tiên đã xuất hiện
trong \(n-1\) vị trí đầu, hoặc \(x_n=1\) và tiền tố trước đó chứa ít nhất
\(k-1\) đầu vào đúng. Vì vậy tổng ít nhất \(k\). Hai cận cho tổng đúng bằng
\(k\). Với mọi vectơ \(X\) có tổng \(k\), gán \(R_{i,j}\) theo tổng tiền tố
thật thỏa hệ thức truy hồi và cả hai điều kiện biên.
\end{proof}

\begin{workedexample}{Exactly-\(3\) trên sáu biến}
Với \(n=6,k=3\), NSC tạo
\[
  \min(1,3)+\min(2,3)+\cdots+\min(5,3)=12
\]
biến thanh ghi, thay vì \(5\cdot3=15\) ô của bảng chữ nhật. Xét
\[
  x=(1,0,1,0,0,1).
\]
Sau năm biến, tổng tiền tố bằng \(2\), nên
\[
  (R_{5,1},R_{5,2},R_{5,3})=(1,1,0).
\]
Điều kiện cuối
\(
R_{5,3}\lor(x_6\land R_{5,2})
\)
đúng chính vì \(x_6=1\) và tiền tố đã đạt ngưỡng \(2\). Nếu đặt thêm
\(x_5=1\), ta có \(R_{5,3}=1\); mệnh đề chống tràn
\(x_6\rightarrow\neg R_{5,3}\) lập tức buộc \(x_6=0\). Hai tình huống này
minh họa lần lượt biên dưới và biên trên của cùng một thanh ghi.
\end{workedexample}

\section{Cơ chế suy diễn lan truyền hai chiều}

Bộ đếm hữu ích vì các cận xuất hiện sớm trong tìm kiếm. Hai tình huống quan
trọng có thể kiểm tra trực tiếp:

\begin{enumerate}
  \item Khi đã có \(k\) đầu vào đúng trong một tiền tố, thanh ghi mức \(k\) của
  tiền tố đó trở thành đúng. Các mệnh đề chống tràn buộc mọi đầu vào còn lại về \(0\).
  \item Khi số vị trí có thể còn nhận giá trị \(1\) vừa đủ để đạt \(k\), điều
  kiện cận dưới và các quan hệ kéo theo lùi lần lượt loại nhánh đặt một vị trí còn
  lại về \(0\).
\end{enumerate}

Hai hướng lan truyền này tương ứng với hai biên ``đã đủ'' và ``phải dùng hết
phần còn lại''. Chúng nên được kiểm tra bằng phép thử đơn vị trên mọi phép gán
bộ phận nhỏ, thay vì suy luận hiệu quả chỉ từ số biến hoặc số mệnh đề.

\section{Dùng một thanh ghi cho nhiều giao diện}

Khi thanh ghi đã có đặc tả ngưỡng, nhiều ràng buộc có thể đọc cùng một mô-đun:
\[
  \sum X\le k
  \quad\Longrightarrow\quad
  \text{các mệnh đề chống tràn},
\]
\[
  \sum X\ge k
  \quad\Longrightarrow\quad
  \text{đầu ra đạt ngưỡng cuối},
\]
\[
  \sum X=k
  \quad\Longrightarrow\quad
  \text{kết hợp hai giao diện trên}.
\]
Nếu cần nhiều cận \(k_1<k_2<\cdots\), bộ đếm nên được tạo đến ngưỡng lớn nhất
rồi cung cấp các đầu ra nhỏ hơn. Cách này tránh xây hai mạng đếm độc lập cho
ràng buộc và hàm mục tiêu.

\section{Nghiên cứu tình huống ngắn: xếp nhóm chơi golf}

\subsection{Từ thiết kế tổ hợp đến mô hình SAT}

\emph{Social Golfer} (bài toán xếp nhóm người chơi golf) vừa là bài toán lập
lịch vừa là bài toán thiết kế tổ hợp. Ba
nguồn khó đan xen với nhau: ràng buộc đúng kích thước nhóm, cấm lặp cặp người
chơi và một nhóm đối xứng lớn do hoán vị tuần, nhóm, người. Vì vậy, hiệu quả
của một phép biểu diễn ràng buộc đếm trên SGP không thể tách khỏi lựa chọn biến
chính và cách phá đối xứng.

Các mô hình SAT trước đó từng dùng thêm chỉ số vị trí trong nhóm, sau đó được
sửa và rút gọn bằng biến người--nhóm--tuần cùng cấu trúc bậc thang
\parencite{triska2012golfer}. Dòng phát triển này minh họa một quy tắc quan
trọng: bỏ một chỉ số không mang nghĩa thường tạo lợi ích lớn hơn tối ưu vài
mệnh đề trong chính bộ biểu diễn ràng buộc đếm. Phép đánh giá NSC vì thế bắt đầu sau
khi phần biểu diễn đối tượng đã được chuẩn hóa.

Một bài toán thử \(g\)-\(p\)-\(w\) có \(n=gp\) người chơi, \(g\) nhóm mỗi tuần,
mỗi nhóm có \(p\) người và \(w\) tuần. Biến ba chỉ số
\[
  G_{i,r,t}=1
  \quad\Longleftrightarrow\quad
  \text{người chơi \(i\) thuộc nhóm \(r\) ở tuần \(t\)}
\]
loại bỏ chỉ số vị trí bên trong nhóm, bởi thứ tự các thành viên trong một nhóm
không mang ý nghĩa.

Mô hình gồm ba họ ràng buộc:
\begin{align*}
  &\sum_{r=1}^{g}G_{i,r,t}=1
    &&\text{với mọi người chơi \(i\), tuần \(t\)},\\
  &\sum_{i=1}^{n}G_{i,r,t}=p
    &&\text{với mọi nhóm \(r\), tuần \(t\)},\\
  &\neg G_{i,r,t}\lor\neg G_{j,r,t}\lor
    \neg G_{i,s,u}\lor\neg G_{j,s,u}
    &&\text{với \(i<j,\ t<u\)}.
\end{align*}
Họ thứ ba cấm một cặp người chơi gặp nhau ở cả nhóm \(r\), tuần \(t\) và nhóm
\(s\), tuần \(u\). NSC được dùng ở họ thứ hai; mỗi cặp nhóm--tuần có một
bộ đếm tam giác riêng \parencite{kieu2026golfer}.

Hai phép chuẩn hóa rẻ được thêm vào mô hình. Tuần đầu được cố định thành các
khối người chơi liên tiếp; ở các tuần sau, một số người chơi neo được gán vào
các nhóm đã định. Chúng loại các hoán vị hiển nhiên của người chơi, nhóm và
tuần mà không thêm biến phụ.

\begin{figure}[tb]
\centering
\begin{tikzpicture}[satfig,x=17mm,y=9mm]
  \node[satlabel] at (0,1.0) {tuần};
  \foreach \g in {1,2,3}
    \node[satlabel] at (\g,1.0) {nhóm \(\g\)};
  \node[satlabel] at (0,0) {\(1\)};
  \node[satlabel] at (0,-1) {\(2\)};
  \foreach \g/\txt in {1/{1,2},2/{3,4},3/{5,6}}{
    \node[satbox,minimum width=15mm,fill=PaleBlue] at (\g,0) {\(\{\txt\}\)};
  }
  \foreach \g/\txt in {1/{1,3},2/{2,5},3/{4,6}}{
    \node[satbox,minimum width=15mm,fill=PaleTeal] at (\g,-1) {\(\{\txt\}\)};
  }
  \draw[decorate,decoration={brace,amplitude=4pt},Gold,thick]
    (3.6,.42)--(3.6,-1.42)
    node[midway,right=5pt,satlabel,text=Gold,align=left]
      {mỗi người đúng một nhóm/tuần;\\mỗi cặp gặp nhiều nhất một lần};
\end{tikzpicture}
\caption{Một lịch khả thi cho bài toán nhỏ \(3\)-\(2\)-\(2\). Mỗi ô nhóm là
một ràng buộc đúng \(2\); tuần đầu có thể được cố định làm lát chuẩn.}
\label{fig:social-golfer-small}
\end{figure}

\begin{resultbox}{Kết quả trên 66 bài toán Social Golfer}
Trong thí nghiệm của \textcite{kieu2026golfer}, cấu hình dùng NSC giải được
56/66 bài toán; các cấu hình ADDER, BDD và CARD cùng giải được 54, còn
BINOMIAL giải được 32. NSC có số lần hết thời hạn thấp nhất; thứ hạng trên từng
bài toán thay đổi theo cấu trúc của công thức. Kết quả làm rõ lợi ích của việc
cân bằng số biến, số mệnh đề và cấu trúc lan truyền.
\end{resultbox}

\subsection{Phân tích kết quả}

Mô hình ba chỉ số loại đối xứng vị trí, còn trạng thái tuần tự tam giác cung
cấp một phép biểu diễn EXK cạnh tranh cho ràng buộc kích thước nhóm. Có thể đo
riêng tác động của NSC bằng cách giữ cố định mô hình SGP, chỉ thay mô-đun kích
thước nhóm, phân tầng theo \(p\), \(g\), \(w\) và tách các bài SAT/UNSAT.

\section{Cài đặt và kiểm thử}

Một bộ sinh NSC nên nhận dãy literal, cận \(k\) và chế độ
\texttt{AMK}/\texttt{ALK}/\texttt{EXK}. Bảng ánh xạ
\((i,j)\mapsto\texttt{varID}\) chỉ chứa ô hợp lệ. Ba lớp kiểm thử là đủ để bắt
phần lớn lỗi biên:

\begin{enumerate}
  \item liệt kê toàn bộ vectơ đầu vào khi \(n\le8\) và so sánh với tổng thật;
  \item với mỗi vectơ hợp lệ, kiểm tra có tồn tại phần mở rộng cho biến \(R\);
  \item với mỗi phép gán một phần, so sánh literal đầu vào mà lan truyền đơn vị
  suy ra với các hệ quả trực tiếp của hai cận.
\end{enumerate}

Các trường hợp \(k=0\), \(k=1\), \(k=n-1\), \(k=n\) cần được kiểm thử riêng. Đây là nơi
quy ước hằng biên và phạm vi chỉ số thường gây lỗi nhất.

\begin{algorithm}[H]
\caption{Sinh NSC cho ràng buộc đúng \(K\)}
\label{alg:nsc}
\begin{minipage}{.94\textwidth}\small
\textbf{Đầu vào:} các literal \(x_1,\ldots,x_n\) và \(0\le k\le n\).\\
\textbf{Đầu ra:} \CNF và ánh xạ ngữ nghĩa của các thanh ghi.
\begin{enumerate}
  \item Xử lý trực tiếp \(k=0\) bằng \(\bigwedge_i\neg x_i\) và \(k=n\)
  bằng \(\bigwedge_i x_i\). Nếu \(k>n/2\), có thể thay
  \(y_i=\neg x_i\) và đếm Exactly-\((n-k)\).
  \item Chỉ cấp phát \(R_{i,j}\) với \(1\le i\le n-1\) và
  \(1\le j\le\min(i,k)\); ghi ánh xạ
  \(R_{i,j}\leftrightarrow[\sum_{q=1}^{i}x_q\ge j]\).
  \item Với từng ô hợp lệ, phân rã tương đương
  \(R_{i,j}\leftrightarrow
    R_{i-1,j}\lor(x_i\land R_{i-1,j-1})\),
  thay \(R_{i,0}=\top\), \(R_{0,j}=\bot\), rồi giản lược các literal hằng.
  \item Sinh các mệnh đề chống tràn của~\eqref{eq:nsc-overflow}.
  \item Sinh điều kiện đạt cận dưới~\eqref{eq:nsc-lower}; trả \CNF và ánh xạ.
\end{enumerate}
\end{minipage}
\end{algorithm}

\section{Miền lựa chọn của NSC}

NSC loại chính xác \(k(k-1)/2\) ô so với bảng chữ nhật, nên tỷ lệ tiết kiệm
biến chỉ là
\[
  \frac{k(k-1)/2}{k(n-1)}
  =\frac{k-1}{2(n-1)}.
\]
Khi \(n\) lớn nhưng \(k\) cố định, tỷ lệ này tiến về \(0\); hằng số của ánh xạ
và các mệnh đề hai chiều có thể quan trọng hơn phần tam giác bị bỏ. Với
\(k\) gần \(n\), đếm các literal bù \(\neg x_i\) đến \(n-k\) thường gọn hơn.
Nếu cùng một tập đầu vào cần nhiều cận hoặc cần thay cận tăng dần, cây tổng và
mạng đếm có đầu ra chia sẻ tự nhiên hơn một NSC chuyên biệt. BDD cũng có thể
thắng khi phép rút gọn hợp nhất được nhiều trạng thái tương đương.

Bộ chọn cho ràng buộc đúng \(K\) dùng \(n\), \(\min(k,n-k)\), số cận sẽ được
truy vấn, nhu cầu giải tăng dần và đặc tả lan truyền làm đặc trưng. NSC phù hợp
với một cận cố định và vùng tam giác bị cắt đủ lớn; cây tổng/mạng đếm phù hợp
khi cần chia sẻ nhiều đầu ra. Kết quả Social Golfer minh họa cách các đặc
trưng này dẫn đến một lựa chọn hiệu quả trên họ bài toán tương ứng.

\section{Bài học thiết kế}

\begin{summarybox}
\begin{enumerate}
  \item Miền trạng thái khả đạt quyết định hình dạng bộ đếm.
  \item Thanh ghi \(R_{i,j}\) có đặc tả ngưỡng rõ ràng, không chỉ là biến
  Tseitin vô danh.
  \item Ràng buộc đúng \(K\) là sự kết hợp của một biên tràn và một điều kiện đạt
  ngưỡng cuối.
  \item Hiệu quả cần được đọc đồng thời qua kích thước, lan truyền và hành vi
  giải trên bài toán khó.
  \item Social Golfer cho thấy lựa chọn biến chính và phép biểu diễn ràng buộc đếm phải
  được thiết kế cùng nhau.
\end{enumerate}
\end{summarybox}

\subsection*{Câu hỏi nghiên cứu}
\begin{enumerate}
  \item Vẽ miền biến \(R_{i,j}\) cho \(n=7,k=3\) và tính phần tiết kiệm so với
  bảng chữ nhật.
  \item Viết các mệnh đề chống tràn cho \(\sum_{i=1}^{6}x_i\le2\).
  \item So sánh hai mô hình Social Golfer có và không có chỉ số vị trí trong
  nhóm về số biến chính và nhóm đối xứng còn lại.
\end{enumerate}
